\documentclass{article}
\usepackage{nips15submit_e,times}
\usepackage{hyperref}
\usepackage{url}

\usepackage[table]{xcolor}
\usepackage{colortbl}
\usepackage{graphicx}
\usepackage{amssymb, amsmath}
\usepackage[normalem]{ulem}


\title{Towards Distributed, Asynchronous Reinforcement Learning on Consumer Compute}


\author{
Gordon Chen\\
School of Computing\\
University of Connecticut\\
\texttt{gbc24001@uconn.edu}\\
}

\newcommand{\fix}{\marginpar{FIX}}
\newcommand{\new}{\marginpar{NEW}}

\nipsfinalcopy

\begin{document}


\maketitle

\begin{abstract}
Deep learning research increasingly depends on access to large amounts of compute. At the same time,
due to recent interest in deep learning, compute is becoming harder to access for student researchers
and hobbyists.
In this paper, we develop a distributed, asynchronous reinforcement learning protocol
that can make use of idle, heterogeneous consumer compute over the internet.
We use IMPALA's V-Trace algorithm to correct for stale policies and show via ablations that it is
necessary in order to make asynchronous training stable.
We introduce trajectory compression and protocol optimizations that allow us to beat a tuned
synchronous baseline, reaching the maximum reward on Pong faster in wall-clock time.
We open-source our code on GitHub at \url{https://github.com/gordonbchen/taming-impala} and hope that
this framework will allow students and researchers to do even more awesome things through
consumer compute pooling.
\end{abstract}


\section{Introduction}
Due to the increased interest in deep learning, compute for student and hobbyist research is becoming harder
to access. Both memory and GPUs are getting more expensive, and cloud compute providers frequently run out of
GPU instances when demand is high. Models are getting bigger and experiments are getting more expensive,
meaning that even relatively small research projects can be constrained by access to compute.

Simultaneously, a large amount of consumer compute, such as desktops and gaming PCs,
sits around unused for most of the day. Many students and hobbyists often have desktops with decent hardware
that spend much of the day idle. Alone, these machines are often not powerful enough
for modern deep learning experiments, even though collectively they may represent a substantial pool
of useful compute.

We believe that reinforcement learning is very well suited for distributed asynchronous training, even
more so that classic supervised learning. In supervised learning, data is collected beforehand and during
training there is really only optimization. On the other hand, reinforcement learning algorithms very
naturally split between alternating rollout and policy / value function optimization phases. RL's unique
property that data is collected online as the model changes makes it ripe for asynchronous training, which
increases update throughput by overlapping rollout and optimization phases at the cost of introducing
the negative effects of stale data.

In this paper, we develop a distributed, asynchronous reinforcement learning protocol that
allows multiple geographically scattered machines to contribute compute to a shared training run
over the internet.
Our protocol follows the general setup from IMPALA, with actor nodes that collect trajectories using local,
often stale, copies of the policy network, and periodically send those trajectories to a central learner
over the internet, and receive updated model parameters in return.

Building upon IMPALA, we introduce protocol optimizations such as rollout compression and a custom message
format to reduce internet bandwidth constraints, allowing for higher update throughput. The flexible protocol we
develop is not strongly tied to any particular neural network architecture or environment. Furthermore,
in contrast to the within-datacenter, distributed training with homogenous compute that IMPALA was meant for,
our protocol is developed for unreliable, bandwidth-limited setups for compute pooling over the internet, and
thus supports actors with heterogeneous compute capabilities (CPU and GPU), as well as resilience to actors
leaving and joining the commpute pool in the middle of a training run.

We thoroughly ablate our protocol, and find that IMPALA's V-Trace is necessary to stabilize asynchronous
training at the levels of policy lag our cross-internet protocol faces.


\section{Background}
In 2018, Espeholt et al. released IMPALA, an asynchronous reinforcment learning architecture that uses
multiple actor nodes to roll out trajectories with a local, possibly stale actor and sends collected
trajectories to a central learner, which optimizes the policy and value function, and broadcasts the updated
model parameters to actors. Critically, the authors introduce V-Trace, an algorithm that uses importance
sampling corrections to adapt policy gradient updates to use rollouts with stale policies. We will now derive
V-Trace.

We first start with the vanilla policy gradient
$$
\nabla_\theta J(\theta) = \mathbb{E}_{\tau \sim \pi_\theta}
\left[ \sum_{t=1}^T \nabla_\theta \log \pi_\theta(a_t|s_t) \ \Phi_t \right]
$$

where $J(\theta)$ is the expected, discounted cumulative reward of a parameterized policy $\pi_\theta$ in some
Markov Decision Process. Rollout trajectories are dented $\tau$, and are made up of states $s_t$, actions $a_t$,
and rewards $r_t$ from times $1..T$.

Here, $\Phi_t$ can represent the cumulative discounted return, Rewards-To-Go, or any general advantage estimate.
For our purposes, we choose $\Phi_t$ to be the baselined $n$-step Bellman target.
\begin{align*}
    \Phi_t &= v_t - V_\phi(s_t) \\
    v_t &= \sum_{t'=t}^{T-1} \gamma^{t'-t} r_t' + \gamma^{T-t} V_\phi(s_T)
\end{align*}

Here, $v_t$ is the $n$-step Bellman target, $V_\phi$ is our parameterized value function baseline,
and $\gamma$ is the future reward discount factor.

The $n$-step Bellman target is essentially the cumulative discounted reward for $n$ steps and the bootstrapped
discounted value of the last state. We can also write the $n$-step Bellman target in terms of TD-error $\delta_t$
as follows:
\begin{align*}
    \delta_t &= r_t + \gamma V_\phi(s_{t+1}) - V_\phi(s_t) \\
    v_t &= V_\phi(s_t) + \sum_{t'=t}^{T-1} \gamma^{t'-t} \delta_t
\end{align*}

Now we can see that the $n$-step Bellman target is the value of the current state, plus discounted TD-error
correction factors.

However, in the context of asynchronous training, we have a problem. In our equation for the $n$-step Bellman
target, $V_\phi(s_t)$ is the expected value of the state $s_t$ if we follow the current policy, but the
rewards and therefore TD-error correction terms are coming from rollouts that were collected on the actor,
so from a stale policy. This mismatch means that $V_\phi(s_t)$ is supposed to be with respect to the current
policy, but will be trained to predict the return coming from stale policies.






% \subsubsection*{Acknowledgments}
% \subsubsection*{References}
\end{document}$
